\section{Experimental Setup}

\subsection{Dataset}

We use the CPJUMP1 pilot~\cite{chandrasekaran2024jump}, a public dataset from the JUMP Cell
Painting Consortium. It contains 303 small-molecule compounds, 160 CRISPR knockouts and 176 ORF
overexpressions measured in U2OS and A549 cells. Most perturbation types were collected at two
timepoints. Each well contains several imaged sites with five fluorescence channels; the three
available brightfield channels are not used in this study.

\subsection{Data Pipeline}

Because the raw images occupy several terabytes, we process one plate at a time. After DINOv3
produces the per-channel features, the raw local copy is removed and the cached features are reused
in every training run. Text prompts are handled similarly: ModernBERT encodes each prompt once, and
the resulting feature is stored for later training.

\subsection{Splits}

We split the data by perturbation rather than by well. A fixed hash assigns each perturbation to
the training, validation or test set in an 80/10/10 ratio, so replicate wells of one perturbation
never appear in different splits. All three perturbation types are represented in each split, and
control wells are excluded from training and retrieval evaluation. The validation set contains
2,220 wells from 98 perturbations, and the test set contains 1,860 wells from 86 perturbations.

\subsection{Retrieval Protocol}
\label{sec:protocol}

Replicate wells share a text description. A split therefore contains $P$ distinct text vectors and
$N$ well vectors. We evaluate four retrieval settings and compute an analytic chance baseline for
each one.

\begin{itemize}
  \item \textbf{Well-level image-to-text.}
        Each of the $N$ wells ranks the $P$ unique texts. One candidate is correct. Random
        Recall@$k$
        is $k/P$.
  \item \textbf{Well-level text-to-image.}
        Each of the $P$ texts ranks all $N$ wells and is scored on the rank of its first correct
        replicate. With $m$ replicate wells among $N$ candidates, random Recall@$k$ is
        $1 - \binom{N-m}{k} / \binom{N}{k}$.
  \item \textbf{Perturbation-level, both directions.}
        The well embeddings of each perturbation are averaged and re-normalized into one profile,
        and the $P$ profiles are ranked against the $P$ texts. This is the unit of the standard
        CPJUMP1 protocol. Random Recall@$k$ is $k/P$.
\end{itemize}
We report Recall@\{1, 5, 10\} and the median rank of the first correct match. Checkpoints are
selected using validation text loss and then evaluated on both validation and test data.

\subsection{Standard CPJUMP1 Benchmark}

We also evaluate exported well embeddings with the standard CPJUMP1 benchmark of Chandrasekaran et
al.~\cite{chandrasekaran2024jump}. We follow the reference pipeline: every well is encoded in fp32,
including controls, and the negative-control mean is removed before scoring. The benchmark uses the
40 plates that meet its standard experimental criteria and covers a short and long timepoint for
each perturbation type. It measures \emph{replicability} between repeated wells, \emph{target
  matching} between perturbations of the same type and \emph{gene--compound matching} across
perturbation types. Results are reported as mean average precision (mAP) and as the fraction of
perturbations that pass the benchmark's significance threshold.

The benchmark's significance test uses random permutations, so fraction retrieved varies slightly
between otherwise identical runs. Replicability mAP is deterministic. We therefore treat small
differences in fraction retrieved as descriptive rather than conclusive.

\subsection{Baselines}

\begin{itemize}
  \item \textbf{CellProfiler}~\cite{carpenter2006cellprofiler} provides the established
        handcrafted-feature baseline. We use its published CPJUMP1 target-matching range because we
        did not rerun the full pipeline.
  \item \textbf{CellCLIP}~\cite{lu2025cellclip} is the closest text-supervised baseline. We report
        its published results and also run the released checkpoint through our short-timepoint
        benchmark pipeline. This comparison is approximate because CellCLIP applies an additional
        KernelPCA correction and the benchmark significance test is stochastic.
  \item \textbf{CWA-MSN}~\cite{huang2025cwamsn} is included as a related batch-aware method. Its
        published benchmarks differ from ours, so a direct numerical comparison is not possible.
\end{itemize}

\subsection{Ablation Design}
\label{sec:ablation_design}

The base model uses binary CWCL labels, random batches, no replicate loss and no plate correction.
The ablations use perturbation-aware batches, a shorter schedule and early stopping. Starting from
this shared control, we add gene-aware soft labels ($\alpha=0.6$), replicate alignment
($\lambda_{\text{rep}}=0.3$) and condition-relative plate offsets one at a time, followed by a run
that combines all three. All runs use seed 42, batch size 256 and the same data split. Since a
single retrieved perturbation changes recall by about one percentage point, small differences
should not be interpreted as reliable effects.

\subsection{Implementation}

All experiments run on one RTX 5080 with 16\,GB of memory. We use AdamW with a learning rate of
$10^{-4}$, weight decay of 0.2, 100 warm-up steps followed by cosine decay, FP16 training and
gradient clipping at 1.0. The CrossChannelFormer has one layer and four attention heads; both
projection heads use a hidden width of 512 and dropout of 0.3. Code, configurations and split
manifests are released with the paper.
